% --- Section: Analysis and Validation ---

\section{Analysis and Validation}
\label{sec:discussion}

\textbf{What the three-ranking divergence means.} The three metrics
are not the same axis under different names. Each has deployment
consequences the others do not capture. Accuracy is the natural
signal for one-shot question answering. For pipelines that
paraphrase prompts or swap inputs (automated tutoring, agentic
workflows), robustness matters more: picking by accuracy alone for
paraphrase-heavy pipelines can be exactly wrong. For systems that
gate decisions on expressed confidence (triage tools, ensemble
selectors), calibration matters, and DeepSeek's ECE of $0.198$
rules it out regardless of accuracy.

\textbf{Reasoning versus pattern matching.} The keyword-vs-judge
gap on assumption-compliance is the clearest evidence for the paper's
title. On the combined base-and-perturbation pool ($n = 2{,}850$),
$20.7\%$ of responses pass-flip between the two scorers. On the
$n=750$ base-only pool, the judge passes $55.7\%$ versus the
keyword rubric's $71.1\%$, a $15.4$-point gap that indicates the
direction: surface-form scoring overcounts substantive
articulation~\cite{du2025icecream}. The calibration analysis
(Table~\ref{tab:calibration-paired}) shows the same shape:
verbalized confidence places Claude and GPT-4.1 in the
well-calibrated tier (ECE $\approx 0.033$), while self-consistency
on three reruns at $T = 0.7$ puts every model past the
$\text{ECE} = 0.6$ ``severe overconfidence'' threshold.

\textbf{Broader implications.} The multi-axis ranking is cheap to
add: the same response corpus feeds robustness once perturbations
exist and calibration once confidence is parsed. Single-number
leaderboards should be reported alongside the multi-axis view, not
in place of it.

\subsection{Decomposing the Divergence: Four Orthogonal Signals}
\label{sec:decomposition-synthesis}

The three-ranking divergence is one coarse-grained view. \S
\ref{sec:extended-diagnostics} decomposes it into four orthogonal
diagnostic signals.
\textit{(i) Model-specific over-credit}
(\S\ref{sec:per-model-disagreement}): Mistral $12.11\times$ and
Claude $7.84\times$ keyword-judge asymmetry, vs GPT-4.1
$2.58\times$ --- surface-token credit is not uniform across the
benchmarked LLMs.
\textit{(ii) Task-specific failure concentration}
(\S\ref{sec:failure-by-task-type}): all 143 L1 failures fall in
Phase~1, with the Bayesian-regression families
(\texttt{REGRESSION}+\texttt{BAYES\_REG}) absorbing $35\%$;
failures hit all five models uniformly within each task type, so
the over-credit problem and the task-difficulty problem are
independent.
\textit{(iii) Axis-specific robustness}
(\S\ref{sec:per-axis-robustness}): rephrasing ($R=0.966$) hurts
more than numerical perturbation ($R=0.978$), with model variance
$3.66\times$ axis variance --- the robustness ranking is a real
model property, but GPT-4.1's $0.9996$ is a near-cancellation of a
rephrasing drop and gains on the other two axes.
\textit{(iv) Phase-specific rank flips}
(\S\ref{sec:per-phase-performance}): Mistral climbs from 5th on
Phase~1 to 2nd on Phase~2 ($+3$ positions), driven by the
assumption-compliance bar dropping on procedural prompts where the
rubric expects no explicit prior or likelihood statement.

A model can be top-ranked on one signal and bottom-ranked on
another --- the four are not a single latent factor under different
names. Mistral concretely illustrates this: 5th on Phase~1 NMACR
but 2nd on Phase~2; the strongest pattern-matcher per
\S\ref{sec:per-model-disagreement} but the second-most-robust per
\S\ref{sec:per-axis-robustness}. Pattern matching as a strategy is
robust --- surface tokens propagate cleanly through paraphrase and
number-swap --- and pays off on tasks where the rubric accepts
surface-correct procedural prose. It collapses only on tasks with
substantive assumption-articulation requirements that the judge
actually checks (Phase~1 conjugate derivations, specifically
Bayesian linear regression).

A single-metric leaderboard does not just collapse the
three-ranking divergence from \S\ref{sec:results}: it collapses
four orthogonal diagnostic axes. The title question admits a more
specific answer than the headline three-ranking captures: pattern
matching is real, model-specific, axis-stable, and rewarded by
procedural task structure. A reader picking a model for Bayesian
work should inspect all four \S\ref{sec:extended-diagnostics}
columns alongside the headline accuracy and calibration numbers ---
the aggregate scoreboard hides which model is genuinely reasoning
and which is fluent at the surface.
